%% * * * * * * * * * * * * * * * * * * * * * * * * * * * * * * * * * * * * * * * * 
%
%   Derivations in High School Physics Chapter Two: Thermodynamics
%
%   Derivation of concepts in thermodynamics relavent to physics competitions and AP courses
%   Author: Justin Ely
%   Email: justinyely@gmail.com
%   
%% * * * * * * * * * * * * * * * * * * * * * * * * * * * * * * * * * * * * * * * * 

\chapter{Thermodynamics}\label{Chapter 2}

%% * * * * * * * * * * * * * * * * 
%
%   2.1: Molecular Thermodynamics
%
%   Includes: 
%
%   Maxwell-Boltzman Distribution
%
%% * * * * * * * * * * * * * * * * 

\section{Statistical Thermodynamics}\label{2.1}

%% * * * * * * * *
%
%   2.1.1: Maxwell-Boltzmann Distribution
%
%% * * * * * * * *
\subsection{\color{Orchid} $\blacktriangleright$ \color{black} The Maxwell-Boltzmann Distribution}\label{2.1.1}
An ideal gas can be represented by a Gaussian distribution of the velocity of the singular particles in the gas, given by the following distribution:
\begin{equation}
    f(v) = \frac{4}{\sqrt{\pi}}\left[\frac{m}{2k_BT}\right]^{3/2}v^2e^{\frac{-mv^2}{2k_BT}}
\end{equation}
\noindent The function is known as the \textbf{Maxwell-Boltzmann Distribution} and although it is a bit messy, it's important to recognize what it means. \\
\\
\noindent The function is a \textbf{Probability Distribution Function (PDF)}, which is a special type of function that describes the way in which a random system is arranged. If we consider an arbitrary PDF $\rho(x)$ where $x$ is some state of the system, $\rho$ has no real meaning at any value of $x$. However, the area under $\rho$ (integral) represents the probability of the sytem taking on a state that lies between the bounds of integration. Thus, any PDF $\rho(x)$ has the property
\begin{equation}
    \int_{-\infty}^{\infty} \rho(x) dx = 1
\end{equation}
\noindent since the system must always be occupying one of its possible states. Similarly, the probability of an event occurring in the interval $(x,x+dx)$ is $\rho(x) dx$. For a container of gas, the particles are all evenly distributed, meaning that the probability of a gas particles having a velocity in the range $(v,v+dv)$ equals the fraction of all the particles with that velocity. Or,
\begin{equation}
    \frac{dN}{N} = f(v) dv
\end{equation}
\noindent where $dN$ is the number of gas particles with a speed in the range $(v, v+dv)$. Recall that the expected value (or mean) of a system is defined as
\begin{equation*}
    \langle x \rangle = \sum_i x_i p(x_i)
\end{equation*}
\noindent where $\langle x\rangle$ is the average of $x$, $x_i$ is an arbitrary state, and $p(x_i)$ is the probability of $x_i$ occurring. It follows that for continuous distributions 
\begin{equation}
    \langle x\rangle = \int_{-\infty}^{\infty} x \rho(x) dx
\end{equation}
\noindent since $x$ is a state, $\rho(x)dx$ is the probability, and we are again summing over all possible states. Similarly, the mean of any function of $x$ is
\begin{equation}
    \langle g(x)\rangle = \int_{-\infty}^\infty g(x)\rho(x) dx.\footnote{The reason why we don't change $\rho(x)$ to say, $g(\rho(x))$ is because the fundamental randomization process remains unchanged, and we are simply rescaling the quantity. For example, if we wanted to find the average of the square of a die roll, we would multiply the outcomes $(1,4,9,16,25,36)$ by the probability $1/6$ and add them up to find the mean, which is the same probability used to calculate the mean value of the die normally. This all makes sense, because we aren't changing anything about the die roll, just the output value of each state. At any rate, these properties will be important when we examine macro-scale properties of systems of small particles.}
\end{equation}
\\
\noindent Going back to the Maxwell-Boltzmann distribution, we can break it into two parts:
\begin{align*}
    f(v) &= \frac{4}{\sqrt{\pi}}\left[\frac{m}{2k_BT}\right]^{3/2}\\
    &\cdot v^2e^{\frac{-mv^2}{2k_BT}}
\end{align*}
\noindent The bottom part of the distribution relates to the actual distribution of the speeds of the particles in the gas\footnote{i.e. it's the part that gives the graph its shape} and is derived using statistical physics and experiment. The top part is a normalization factor, which gives the function the property from eq. 2.2. The Maxwell-Boltzmann distribution is special, though, since it's normalized in such a way that the area under the curve from $0$ to $\infty$ is one, instead of the area along the entire $x$ axis. To see why, consider the following graph of $f(v)$:
\begin{figure}[h]
    \centering
    \begin{asy}
        size(8cm);
        
        // Function
        real f(real x) {
          return 3*x^2 * exp(-x^2);
        }
        
        // Axes
        draw((-3.5, 0)--(3.5, 0), Arrow); // x axis
        draw((0, -0.1)--(0, 2), Arrow); // y axis
        
        // Graph
        draw(graph(f, -3, 3), blue+linewidth(1));
        
        // Labels
        label("$v$", (3.5, 0), SE);
        label("$f(v)$", (0, 2), NW);
        //label("$f(v)$", (1.5, 0.2), NE);
    \end{asy}
    \caption{Maxwell-Boltzmann Distribution}
    \label{2.1}
\end{figure}

\noindent Since objects cannot have negative speed, all of the particles in a gas must lie on the right side of the graph, meaning 
\begin{equation}
    \int_0^\infty f(v)dv = 1.
\end{equation}
\noindent This somewhat unusual normalization will be important in \ref{2.1.2} when we integrate over all particles $\bigstar$

%% * * * * * * * *
%
%   2.1.2: RMS Velocity
%
%% * * * * * * * *
\subsection{\color{Orchid} $\blacktriangleright$ \color{black} Root Mean Square Velocity}\label{2.1.2}
The root mean square (RMS) velocity is defined as the square root of the average of the squares of the velocities of the particles in a system. Mathematically (and possibly much more simply),
\begin{equation}
    v_{RMS} = \sqrt{\langle v^2\rangle}.
\end{equation}
\noindent RMS quantities appear a lot in statistical physics, and they are used preferential to normal averages for a couple reasons: \\
\\
\noindent 1: The average value of an oscillating function may be zero, independent of the amplitude of the oscillations. However, it's hardly intuitive for a wildly-oscillating function to be described the same way as a barely-changing function. The RMS value fixes this discrepancy since it squares all values before taking an average, providing a more accurate representation of a quantity.\\
\\
\noindent 2: As we see in \ref{2.1.3}, since kinetic energy is defined with the square of the velocity, average velocity is useless if we want to calculate the average kinetic energy. Thus, $v_{RMS}$ is used instead of the normal average.\\
\\
\noindent In any case, since we know the PDF for particles in a gas, we can use the definitions listed in \ref{2.1.1} to find a general expression. First, we can find the quantity $\langle v^2\rangle$ by (notice the bounds of integration from eq. 2.6)
\begin{align}
    \langle v^2\rangle = \int_{0}^{\infty}v^2f(v)dv
\end{align}
\noindent where $f(v)$ is the Maxwell-Boltzmann distribution. Plugging in $f(v)$ and rearranging yields
\begin{equation}
    \frac{2}{\sqrt{\pi}} a^{3/2}\int_{0}^\infty v^4e^{-av^2}dv.\nonumber
\end{equation}
\noindent Where we use $a=m/2k_BT$. Integrating by parts with $u=v^3,du = 3v^2dv$ and $dw = ve^{-av^2}dv, w=-\frac{1}{2a}e^{-av^2}$, we have
\begin{equation}
    \frac{2}{\sqrt{\pi}} a^{3/2}\left[\left.-\frac{v^3}{2a}e^{-av^2}\right|_{-\infty}^\infty + \frac{3}{2a}\int_{0}^\infty v^2e^{-av^2}dv\right]
\end{equation}
\noindent The first term in the brackets is zero since the exponential decays to zero at infinity. To find the integral on the right, we could use integration by parts again, but instead we can use Feynman's trick for a quicker solution. By properties of the Gaussian integral \footnote{I won't prove it here, but the derivation is very interesting and you should definitely check it out!},
\begin{equation*}
    \int_{0}^\infty e^{-av^2}dv =\frac12 \sqrt{\frac{\pi}{a}}.
\end{equation*}
\noindent Comparing this property with the integral in eq. 2.8, we can observe the following clever relation
\begin{equation*}
    \int_{0}^\infty v^2e^{-av^2}dv = -\frac{d}{da}\int_{0}^\infty e^{-av^2}dv
\end{equation*}
\noindent where we take a derivative with respect to $a$, since the integral is a function of $a$ \footnote{$v$ is a dummy variable, and is integrated out}. Thus, the expression in eq. 2.8 becomes
\begin{align}
    \langle v^2\rangle &= \frac{4}{\sqrt{\pi}} a^{3/2}\left[\frac{3}{2a}\left(-\frac{d}{da} \frac12\sqrt{\frac{\pi}{a}}\right)\right] \7
    &= \frac{4}{\sqrt{\pi}} a^{3/2}\left[\frac{3}{2a}\left(\frac14\frac{\sqrt{\pi}}{a^{3/2}}\right)\right] \7
    &= \frac{3}{2a}.
\end{align}
\noindent The RMS speed is the square root of this value, or
\begin{align}
    v_{rms} &= \sqrt{\langle v^2\rangle} \7
    &= \sqrt{\frac{3}{2a}} \7
    &= \sqrt{\frac{3k_BT}{m}}
\end{align}
$\bigstar$

%% * * * * * * * *
%
%   2.1.3: Internal Energy
%
%% * * * * * * * *
\subsection{\color{RoyalBlue} $\blacktriangleright$ \color{Orchid} $\blacktriangleright$ \color{black} Internal Energy}\label{2.1.3}
The total internal energy of a thermodynamic system is defined as the total energy contained within it. For an ideal gas, since there are no internal potentials, this energy is simply the total kinetic energy of the gas particles, or
\begin{equation}
    U_{int} = N\langle K\rangle
\end{equation}
\noindent where $N$ is the number of particles and $\langle K\rangle$ is the number of gas particles. Using the mechanical definition of kinetic energy,
\begin{equation*}
    \langle K\rangle = \frac12 m\langle v^2\rangle = \frac12 mv_{rms}^2
\end{equation*}
\noindent Hence, the total internal energy of a monatomic gas is 
\begin{equation}
    {U_{int} = \frac32K_BT}.
\end{equation}
\noindent Alternatively, if we define $n$ to be the number of moles of gas constant and the constant $R$ to be Avogadro's number times $k_B$, we have
\begin{equation}
    U_{int} = \frac32 (nN_A)k_BT = \frac32 nRT
\end{equation}
\noindent Notice how this result is specifically for \textit{monatomic} gas, and you can think of the equation as the sum of the average kinetic energy in each degree of freedom (x,y,z axis), which gives rise to the factor of $3/2$. 
\begin{equation*}
    K_i = \frac12mv_x^2 + \frac12mv_y^2 + \frac12mv_z^2
\end{equation*}
\noindent For diatomic gasses, there are two\footnote{Two, not three, since there is symmetry about two axes. For objects with less symmetry, there will be three rotational axes (Try it! Try to find more than two unique rotational axes for a pencil, then try the same with a smartphone!).} additional rotational degrees of freedom, namely
\begin{equation*}
    K_i = \frac12mv_x^2 + \frac12mv_y^2 + \frac12mv_z^2 + \frac12I_x\omeg_x^2 + \frac12I_y\omeg_y^2
\end{equation*}
\noindent Logically, the total internal energy should be
\begin{equation}
    U_{int} = \frac52 Nk_BT
\end{equation}
$\bigstar$

%% * * * * * * * *
%
%   2.1.4: Ideal Gas Law
%
%% * * * * * * * *
\subsection{\color{Orchid} $\blacktriangleright$ \color{black} Ideal Gas Law}\label{2.1.4}
The famous Ideal Gas Law $PV = nRT$ was originally derived by experiment as an amalgamation of other emperical laws, namely Boyle's Law, Charles' Law, Gay-Lussac's Law, and Avogadro's Law. However, using the kinetic theory of gasses, it's possible to derive the law from first principles.\\
\\
We start by considering a cube with side length $L$ filled with a monatomic ideal gas. Since the gas exerts equal pressure in all directions, if we can relate the pressure of particles bouncing off of a single wall to the box's volume and temperature, we simply get the ideal gas law. For this derivation, we look at the particles bouncing off of the wall in the positive x axis (shown in blue). In a small time $\Delta t$, some particles traveling in the positive $x$ direction collide with the blue surface, and those particles take up the volume defined by $\Delta l = \overline{v_x}\Delta t$ where $\overline{v_x}$ is the average $x$ component of the velocities of the gas particles.
\begin{figure}[h]
    \centering
    \begin{asy}
        size(6cm);

        //Verticies
        triple B = (1, 0, 0); // Bottom-right-front
        triple C = (1, 1, 0); // Top-right-front
        triple D = (0, 1, 0); // Top-left-front
        triple E = (0, 0, 1); // Bottom-left-back
        triple F = (1, 0, 1); // Bottom-right-back
        triple G = (1, 1, 1); // Top-right-back
        triple H = (0, 1, 1); // Top-left-back

        triple s1 = (1, 0.7, 0);
        triple s2 = (1, 0.7, 1);
        triple s3 = (0, 0.7, 1);

        triple a0 = (0, 1.3, 0);
        triple ax = (0, 1.6, 0);
        triple ay = (0.3, 1.3, 0);
        triple az = (0, 1.3, 0.3);

        //Labels
        Label A = Label("$A$", position = MidPoint, align=NE+NE+NE);
        Label vt = Label(scale(0.8)*rotate(-9)*"$\Delta l$", position = MidPoint, align = NE);
        label(scale(0.8)*"$x$", ax, E);
        label(scale(0.8)*"$y$", ay, SW);
        label(scale(0.8)*"$z$", az, N);

        //Face
        draw(surface(C--G--H--D--cycle), surfacepen=RGB(115, 185, 250));
        
        //Edges
        draw(s1--s2, black);
        draw(B--C, black);
        draw(C--D, black);
        draw(s2--s3, black);
        
        draw(E--F, black);
        draw(F--G, black);
        draw(G--H, black);
        draw(H--E, black);
        
        draw(s2--G, black, L=vt);
        draw(B--F, black);
        draw(C--G, black);
        draw(D--H, black);

        draw(a0--ax, Arrow3(black));
        draw(a0--ay, Arrow3(black));
        draw(a0--az, Arrow3(black));
        
        // Set 3D view
        currentprojection = perspective(4, 2, 2);
    \end{asy}
    \caption{Box with gas particles randomly moving inside. In a certain time $\Delta t$, some particles inside the section $\Delta l$ hit the blue wall.}
    \label{2.2}
\end{figure}

\noindent In this section of the cube, the number of particles is the number density of particles in the entire box multiplied by the volume of the section, since the particles are all evenly distributed. That number is 
\begin{equation}
    \frac{N}{V}\Delta lA =\frac{NAv_x\Delta t}{V}. \nn
\end{equation}
\noindent Of these particles that move on the $x$ axis, half of them move towards the blue surface and half away. However, only those that move toward the surface exert any pressure, bringing our total considered particles to 
\begin{equation}
    \frac{NAv_x\Delta t}{2V}. 
\end{equation}
\noindent Since the gas is ideal, it collides elastically withe the walls of the container, imparting a force of $F = 2mv_x/\Delta t$ each. Thus, using eq. 2.16, the total force imparted on the blue side is 
\begin{equation}
    \frac{NAmv_x^2}{V}.
\end{equation}
\noindent Since $P = F/A$, the average pressure on the blue side is 
\begin{equation}
    \langle P\rangle = \frac{Nm\langle v_x^2\rangle}{V}.
\end{equation}
\noindent Since the particles are moving in all directions randomly, we can write the relation $v_{rms}^2 = \langle v^2\rangle = \langle v_x^2\rangle + \langle v_y^2\rangle + \langle v_z^2\rangle = 3\langle v_x^2\rangle$, which implies that $\langle v_x^2\rangle = \frac{k_BT}{m}$. Substituting everything in, we have\footnote{I removed the angle brackets around $P$ here since the $P$ in $PV = nRT$ represents the average pressure of the gas in a certain state of the gas. More specifically, the $P$ in $PV = nRT$ represents the "average" over all of the microstates that a given macrostate encompasses.}
\begin{align}
    P &= \frac{Nk_BT}{V} \7
    PV &= Nk_BT \\
    PV &= nRT
\end{align}
$\bigstar$

\clearpage

%% * * * * * * * * * * * * * * * * 
%
%   2.2: Macroscopic Thermodynamics
%
%   Includes: 
%
%   Work Done on an Ideal Gas
%
%% * * * * * * * * * * * * * * * * 

\section{Macroscopic Thermodynamics}\label{2.2}

%% * * * * * * * *
%
%   2.2.1 Work Done on Ideal Gas
%
%% * * * * * * * *
\subsection{\color{RoyalBlue} $\blacktriangleright$ \color{Orchid} $\blacktriangleright$ \color{black} Work Done on an Ideal Gas} \label{2.2.1}
Finding the work done on an ideal gas is important, since we can combine it with the first law of thermodynamics to derive some interesting relations (as seen later). To start, consider cylinder of gas of volume $V$ topped with a movable piston of area $A$. If we apply a force $F$ to the piston and depress it by a \textit{distance} $\Delta x$, the work done by this force is simply
\begin{equation*}
    W = F\Delta x.
\end{equation*}
\noindent Dividing and multiplying by $A$ yields
\begin{align}
    W &= \frac{F}{A}A\Delta x \7
    W &= -P\Delta V.
\end{align}
\noindent Notice here that $A\Delta x = -\Delta V$. Since the cylinder is compressed, $\Delta V$ is negative. Thus, we need to add a negative sign to make things work out. This negative sign makes sense, since we'd expect that compressing a gas would do positive work. Why? Similar to a spring, when you compress the gas you work against a force, you're increasing the energy of the system. At a molecular level, you are simply increasing the momentum of the particles in the box as they collide with the surface of the piston! (see \ref{2.1.4}) $\bigstar$

%% * * * * * * * *
%
%   2.2.2. Heat Capacity of Gasses and Meyer's Relation
%
%% * * * * * * * *
\subsection{\color{RoyalBlue} $\blacktriangleright$ \color{Orchid} $\blacktriangleright$ \color{black} Heat Capacity of Gasses and Meyer's Relation}\label{2.2.2}
The specific heat capacity $C$ of a substance is defined as the ratio of the amount of heat it takes to raise the temperature of a substance per mole of substance. In other words,
\begin{equation}
    Q = nC\Delta T
\end{equation}
\noindent Solids and liquids have relatively constant heat capacities over most temperature intervals, owing to their low thermal expansion. Gasses, on the other hand, have variable heat capacities, primarily because of their ability to change volume. Thus, it's helpful to define a few heat capacities for gasses, and we can even derive some useful relations between them.\\
\\
\noindent First is the specific heat capacity at constant volume or $C_v$. Since the volume of the gas is constant, no work can be done on it and thus it can only change energy via heat transfer. Plugging in appropriate expressions into the First Law of Thermodynamics, we have (for a monatomic gas)
\begin{align}
    \Delta U &= Q \7
    \frac32 nR\Delta T &= nC_v\Delta T \7
    C_v &= \frac32R
\end{align}
\noindent The other notable heat capacity is the heat capacity at constant pressure. Writing out the First Law of Thermodynamics for a constant pressure pressure yields
\begin{align}
    \Delta U &= W +Q \7
    \frac32 nR\Delta T &= -P\Delta V + nC_p\Delta T .
\end{align}
\noindent Since pressure is constant, we have from the Ideal Gas Law that 
\begin{equation}
    \Delta V = \frac{nR\Delta T}{P}
\end{equation}
\noindent Substituting eq. 2.25 into 2.24 and simplifying gives us
\begin{align}
    \frac32 nR\Delta T &= -nR\Delta T + nC_p\Delta T \7
    C_p &= \frac52R \7
    C_p &= C_v + R.
\end{align}
\noindent Equation 2.26 is known as Mayer's Relation and provides a convenient relation between the heat capacity of a gas at constant volume and pressure. An important thing to see is that it doesn't say anything about the values of $C_v$ or $C_p$, just their relation. In fact, for diatomic gasses, the values for $C_v$ and $C_p$ are $5R/2$ and $7R/2$ respectively, which is a consequence of the phenomenon discussed towards the end of \ref{2.1.3} $\bigstar$ 

%% * * * * * * * *
%
%   2.2.3: Adiabatic Curve
%
%% * * * * * * * *
\subsection{\color{Orchid} $\blacktriangleright$ \color{black} The Adiabatic Curve} \label{2.2.3}
On a $P/V$ diagram, isobaric processes are a horizontal line, isochoric processes are a vertical line, and isothermal processes are a $1/x$ curve. These are all pretty intuitive, and can be derived in a few steps. Less intuitive is the curve for an adiabatic process, which we will derive below. By definition of an adiabatic process, any changes in internal energy of a system must arise only from work done on the system, or
\begin{equation}
    dU = \frac32nRdT = \delta W
\end{equation}
\noindent where $\delta$ represents the a small change in a path dependent quantity. Using \ref{2.2.1} and the Ideal Gas Law, we have
\begin{align}
    \frac32nRdT &= -PdV \7
    nC_vdT &= -nR\frac{TdV}{V} \7
    \frac{C_v}{R}\int\frac{dT}{T} &= -\int\frac{dV}{V} \7
    \frac{C_v}{R}\ln T &= -\ln V + C
\end{align}
\noindent where $C$ is some constant. Rewriting $R$ in terms of $C_p$ and $C_v$ as well as defining the ratio $\gamma = C_p/C_v$, we have
\begin{align}
    \ln T &= \frac{C_v-C_p}{C_v}\ln V + C \7
    \ln T &= (1-\gamma)\ln V +C \7
    TV^{\gamma-1} &= C
\end{align}
\noindent which itself is a perfectly good relation between quantities for an adiabatic process. However, since we're graphing things on a $P/V$ diagram, it's nice to have a relation between $P$ and $V$. Substituting the temperature using the Ideal Gas Law,
\begin{align}
    \frac{PV}{nR}V^{\gamma-1} &= C \7
    PV^\gamma &= const.
\end{align}
\noindent Since $\gamma>1$, this curve looks similar to the curve for isothermal processes, but is a bit steeper $\bigstar$

%% * * * * * * * *
%
%   2.2.4: Carnot's Theorem
%
%% * * * * * * * *
\subsection{\color{Orchid} $\blacktriangleright$ \color{black} Carnot's Theorem} \label{2.2.4}
Carnot's theorem is a direct result of the second law of thermodynamics, and is an important result for engine design. It states that the most efficient heat engine that runs between two heat reservoirs is a perfectly reversible engine. The corollary of this statement is that all reversible engines running between two reservoirs of different temperature have the same efficiency. \\
\\
\noindent We start the proof by considering a reversible heat engine $R$ with efficiency $\eta_R$ and an arbitrary heat engine $E$ with efficiency $\eta_E$ operating between a hot and cold reservoir.

\begin{figure} [h]
    \centering
    \begin{asy}
        size(7cm);
        
        // Colors
        pen hotColor = red;
        pen coldColor = blue;
        pen workPen = darkgreen;
        
        // Reservoirs
        path hotReservoir = (0,6)--(8,6)--(8,7)--(0,7)--cycle;
        path coldReservoir = (0,0)--(8,0)--(8,1)--(0,1)--cycle;
        
        fill(hotReservoir, hotColor + opacity(0.3));
        label("Hot Reservoir ($T_H$)", (4,6.5), align=Center);
        
        fill(coldReservoir, coldColor + opacity(0.3));
        label("Cold Reservoir ($T_C$)", (4,0.5), align=Center);
        
        // Engines
        pair R_pos = (2,3.5);
        pair I_pos = (6,3.5);
        
        draw(circle(R_pos, 0.5));
        label("R", R_pos);
        draw(circle(I_pos, 0.5));
        label("E", I_pos);
        
        // Heat transfer arrows
        draw(R_pos + (0,2.5) -- R_pos + (0,0.5), L=Label("$Q_{h}$", position=MidPoint, align=E), arrow=Arrow(TeXHead), hotColor);
        draw(R_pos + (0,-0.5) -- R_pos + (0,-2.5), L=Label("$Q_{c}$", position=MidPoint, align=E), arrow=Arrow(TeXHead), coldColor);
        
        draw(I_pos + (0,2.5) -- I_pos + (0,0.5), L=Label("$Q_{h}'$", position=MidPoint, align=E), arrow=Arrow(TeXHead), hotColor);
        draw(I_pos + (0,-0.5) -- I_pos + (0,-2.5), L=Label("$Q_{c}'$", position=MidPoint, align=E), arrow=Arrow(TeXHead), coldColor);
        
        // Work arrows
        draw(R_pos + (0.5,0) -- R_pos + (1.5,0), L=Label("$W$", position=EndPoint, align=E), arrow=Arrow(TeXHead), workPen);
        draw(I_pos + (0.5,0) -- I_pos + (1.5,0), L=Label("$W$", position=EndPoint, align=E), arrow=Arrow(TeXHead), workPen);
    \end{asy}
    \caption{A reversible and irreversible engine running between a hot and cold reservoir.}
    \label{2.3}
\end{figure}

\noindent Both engines output work $W$ but $R$ takes in $Q_h$ and outputs $Q_c$ to $T_C$ while $E$ takes in $Q_h'$ and outputs $Q_c'$ to $T_C$. For the sake of the proof, assume that $\eta_E > \eta_R$, i.e.
\begin{gather}
    \frac{W}{Q_h'} > \frac{W}{Q_h} \7
    Q_h >Q_h'.
\end{gather}
\noindent Now, consider running $R$ in reverse using the work done by $E$, so that $R$ extracts heat $Q_c$ from $T_C$ and deposits heat $Q_h$ into $T_H$.

\begin{figure} [h]
    \centering
    \begin{asy}
        size(7cm);
        
        // Colors
        pen hotColor = red;
        pen coldColor = blue;
        pen workPen = darkgreen;
        
        // Reservoirs
        path hotReservoir = (0,6)--(8,6)--(8,7)--(0,7)--cycle;
        path coldReservoir = (0,0)--(8,0)--(8,1)--(0,1)--cycle;
        
        fill(hotReservoir, hotColor + opacity(0.3));
        label("Hot Reservoir ($T_H$)", (4,6.5), align=Center);
        
        fill(coldReservoir, coldColor + opacity(0.3));
        label("Cold Reservoir ($T_C$)", (4,0.5), align=Center);
        
        // Engines
        pair R_pos = (2,3.5);
        pair I_pos = (6,3.5);
        
        draw(circle(R_pos, 0.5));
        label("R", R_pos);
        draw(circle(I_pos, 0.5));
        label("E", I_pos);
        
        // Heat transfer arrows
        draw(R_pos + (0,0.5) -- R_pos + (0,2.5), L=Label("$Q_{h}$", position=MidPoint, align=E), arrow=Arrow(TeXHead), hotColor);
        draw(R_pos + (0,-2.5) -- R_pos + (0,-0.5), L=Label("$Q_{c}$", position=MidPoint, align=E), arrow=Arrow(TeXHead), coldColor);
        
        draw(I_pos + (0,2.5) -- I_pos + (0,0.5), L=Label("$Q_{h}'$", position=MidPoint, align=E), arrow=Arrow(TeXHead), hotColor);
        draw(I_pos + (0,-0.5) -- I_pos + (0,-2.5), L=Label("$Q_{c}'$", position=MidPoint, align=E), arrow=Arrow(TeXHead), coldColor);
        
        // Work arrows
        //draw(R_pos + (0.5,0) -- R_pos + (1.5,0), L=Label("$W$", position=EndPoint, align=E), arrow=Arrow(TeXHead), workPen);
        draw(I_pos + (-0.5,0) -- R_pos + (0.5,0), L=Label("$W$", position=MidPoint, align=N), arrow=Arrow(TeXHead), workPen);
    \end{asy}
    \caption{A modified setup.}
    \label{2.4}
\end{figure}

\noindent Analyzing the entire system at this state, the two engines extract heat $Q_c - Q_c' > 0$\footnote{$Q_h = W + Q_c$ $\Rightarrow$ e.q. 1.31 $\Rightarrow$ $Q_c > Q_c'$} from $T_c$ and deposit heat $Q_h - Qh' > 0$ into $T_h$. Without any external work done on the system, this transfer of heat from a cold reservoir to a hot reservoir is a direct violation of the second law of thermodynamics, meaning that our initial assumption that $\eta_E > \eta_R$ must be false. Hence, there exists no engine more efficient than a reversible engine operating between two reservoirs. \\
\\
If reversible engines are the most efficient engine between to temperature reservoirs, they must have the same efficiency, which is the maximum possible efficiency. Thus, the corollary mentioned earlier is also true. In real life applications, Carnot's theorem gives the upper bound on efficiency of a given heat engine, although most commercial engines don't come close to the ideal Carnot engine $\bigstar$


%% * * * * * * * *
%
%   2.2.5: Carnot Engine Efficiency
%
%% * * * * * * * *
\subsection{\color{Orchid} $\blacktriangleright$ \color{black} Efficiency of a Carnot Engine}\label{2.2.5}
The Carnot Engine is a perfectly reversible engine operating between a hot and a cold temperature. As shown in \ref{2.2.4}, the efficiency of any such engine is the same for two reservoirs of given temperatures and is the upper bound for efficiency of an engine operating between the temperatures. A common way to construct the Carnot Engine is as follows (and as depicted in the diagram):
\begin{enumerate}
    \item An isothermal expansion at hot temperature $T_H$
    \item An adiabatic expansion 
    \item An isothermal compression at cold temperature $T_C$
    \item An adiabatic compression
\end{enumerate}

\begin{figure}[h]
    \centering
    \begin{asy}
        size(8cm,6cm,IgnoreAspect);
        
        draw((0,0) -- (5,0), arrow=Arrow(TeXHead), 1+black, L=Label("$V$", position=MidPoint, align=S));
        draw((0.5,-.5) -- (0.5,5), arrow = Arrow(TeXHead), 1+black, L=Label("$P$", position=MidPoint, align=W));
        
        // Isothermal processes
        path isothermal_expansion = graph(new real(real V){return 4/V;},1,2);
        path isothermal_compression = graph(new real(real V){return 2/V;},4,2);
        
        // Adiabatic processes
        path adiabatic_expansion = graph(new real(real V){return 8/(V^2);},2,4);
        path adiabatic_compression = graph(new real(real V){return 4/(V^2);},2,1);
        
        // Draw all paths in black
        draw(isothermal_expansion, 0.5+red, Arrow(Relative(0.5)));
        draw(isothermal_compression, 0.5+blue, Arrow(Relative(0.5)));
        draw(adiabatic_expansion, 0.5+black, Arrow(Relative(0.5)));
        draw(adiabatic_compression, 0.5+black, Arrow(Relative(0.5)));
        
        // Process labels
        label("1", (1.5,2.67), black, align=NE);
        label("2", (2.8,1.1), black, align=NE);
        label("3", (2.9,0.55), black, align=SE);
        label("4", (1.3,1.90), black, align=W);
        
        // State labels
        label("$A$",(1,4),NW);
        label("$B$",(2,2),NE);
        label("$C$",(4,0.5),SE);
        label("$D$",(2,1),SW);

        // Dots
        dot((1,4), 4+black);
        dot((2,2), 4+black);
        dot((4,0.5), 4+black);
        dot((2,1), 4+black);
    \end{asy}
    \caption{Carnot Cycle, with two adiabatic and two isothermal processes. The hot temperature $T_H$ is depicted in red while the cold temperature $T_C$ is depicted in blue.}
    \label{2.5}
\end{figure}

\noindent To find the efficiency of the engine, consider first the definition of engine efficiency, which is the ratio of the work output and the heat input. We can thus break each process down and see where heat is being inputted and where work is extracted (note that all of the following values are magnitudes).
\begin{enumerate}
    \item Heat $Q_H$ is inputted and work $W_1$ is extracted. Since $\Delta U =0$, $Q_H = W_1$.\footnote{Note here that the volume increases at a positive pressure, meaning that the quantity $W = -P\Delta V$ is negative, but the work that the gas does is positive. Why? It's because the quantity $-P\Delta V$ is the work done \textit{on} the gas, not \textit{by} the gas, i.e. it gives the work inputted into the gas, not the work extracted. An intuitive way of thinking of it is like this: it takes effort to compress a piston and decrease its volume (positive work input) and if released from a compressed state, the piston springs outwards (positive work output).}
    \item No heat is exchanged but work $W_2$ is extracted.\footnote{Again, $W_2$ is a positive value despite there being an increase in volume.}
    \item Heat $Q_C$ is extracted and work $W_3$ is inputted. Since $\Delta U=0$, $Q_C = W_3$.
    \item No heat is exchanged but work $W_4$ is inputted.
\end{enumerate}
\noindent By conservation of energy, the work output of the engine is $Q_H - Q_C$, meaning that the efficiency of the engine is 
\begin{equation}
    \epsilon_{carnot} = \frac{W_{out}}{Q_H} = \frac{Q_H-Q_C}{Q_H}.
\end{equation}
% \noindent As mentioned above, $W_1 = Q_H$, so eq. 1.32 reduces to 
% \begin{equation}
%     \epsilon_{carnot} = 1 + \frac{W_2}{Q_H}.
% \end{equation}
% \noindent Since process $2$ takes the engine from $T_H$ to $T_C$ by outputting work, $W_2$ is simply the difference in energies between $B$ and $C$.
% \begin{equation}
%     W_2 = nR(T_H-T_C)
% \end{equation}
\noindent Since $Q_H = W_1$, we can find $Q_H$ by solving for the area under process $1$:
\begin{align}
    Q_H = W_1 &= \int_{V_A}^{V_B}PdV \7
    &= nRT_H\int_{V_A}^{V_B}\frac{dV}{V} \7
    &= nRT_H\ln\left(\frac{V_B}{V_A}\right)
\end{align}
\noindent where $V_A$ and $V_B$ are the volumes of the gas at state $A$ and $B$, respectively. By the same logic, 
\begin{equation}
    Q_C = W_3 = \int_{V_D}^{V_C}PdV = nRT_C\ln\left(\frac{V_C}{V_D}\right)
\end{equation}
\noindent Plugging eq. 2.33 and 2.34 into 2.32, we have
\begin{equation}
    \epsilon_{carnot} = 1-\frac{T_C}{T_H}\cdot\frac{\ln(V_B/V_A)}{\ln(V_C/V_D)}
\end{equation}
\noindent From the derivation in \ref{2.2.3} (specifically eq. 2.29), we can relate the volumes between the start and end points of the adiabatic processes
\begin{align*}
    \frac{T_H}{T_C} = \left(\frac{V_C}{V_B}\right)^{\gamma-1} &= \left(\frac{V_D}{V_A}\right)^{\gamma-1} \7
    \frac{V_C}{V_D} &= \frac{V_B}{V_A}.
\end{align*}
\noindent With this result, eq. 2.35 simply reduces to 
\begin{equation}
    \boxed{\epsilon_{carnot} = 1-\frac{T_C}{T_H}}.
\end{equation}
$\bigstar$